\subsection{Factoriales Precalculados para Coeficientes Binomiales}

\graybold{Preguntas guía:}

\begin{itemize}
\item ¿Necesito muchos coeficientes binomiales $\binom{a}{b}$ módulo un primo, con $a,b$ acotados por un límite fijo razonable ($\le 10^6$-$10^7$)?
\item ¿Calcular cada factorial desde cero por consulta sería demasiado lento dado el volumen de consultas?
\item ¿El módulo es primo, permitiendo invertir factoriales con el Pequeño Teorema de Fermat?
\end{itemize}

\graybold{¿Para qué sirve?} Responder en \bigO{1} cualquier consulta de coeficiente binomial módulo un primo, tras una única precomputación de factoriales e inversos de factoriales. Es la herramienta base de casi cualquier problema de combinatoria modular.

\graybold{Complejidad:} \bigO{MAXN + log p} para precomputar; \bigO{1} por consulta.

\graybold{Fundamento teórico:} $\binom{a}{b} = \frac{a!}{b!(a-b)!}$. Se precalculan todos los factoriales módulo $p$ hasta el máximo valor posible con una multiplicación acumulada, en \bigO{MAXN}. Para los inversos, en vez de invertir cada factorial por separado (\bigO{MAXN \log p} con Fermat repetido), se invierte SOLO el último factorial con Fermat (\bigO{\log p}), y el resto se derivan hacia atrás con la identidad $\text{invfact}[i-1] = \text{invfact}[i] \cdot i \bmod p$ (ya que $\text{fact}[i] = \text{fact}[i-1] \cdot i$, invertir ambos lados da esa relación), quedando toda la tabla en \bigO{MAXN} adicional.

\graybold{Ejercicio aplicado}

Dados $n$ pares $(a,b)$, calcular $\binom{a}{b} \bmod (10^9+7)$ para cada uno.

\graybold{I/O:} n ≤ 10\textsuperscript{5}, a ≤ 10\textsuperscript{6} → lectura total con tokenización (patrón 2); la precomputación se hace una sola vez para todas las consultas.

\graybold{Código solución}

\begin{lstlisting}[language=Python]
import sys

MOD = 1000000007
MAXN = 1000001

def solve():
    fact = [1] * MAXN
    for i in range(1, MAXN):
        fact[i] = fact[i - 1] * i % MOD

    inv_fact = [1] * MAXN
    inv_fact[MAXN - 1] = pow(fact[MAXN - 1], MOD - 2, MOD)   # Fermat: unico inverso modular necesario
    for i in range(MAXN - 1, 0, -1):
        inv_fact[i - 1] = inv_fact[i] * i % MOD

    data = sys.stdin.buffer.read().split()
    n = int(data[0])
    idx = 1
    out = []
    for _ in range(n):
        a = int(data[idx]); b = int(data[idx + 1]); idx += 2
        out.append(str(fact[a] * inv_fact[b] % MOD * inv_fact[a - b] % MOD))
    sys.stdout.write("\n".join(out) + "\n")

solve()
\end{lstlisting}